\chapter{Task A. Directed Link Prediction on the Synthetic Pharmacotherapy Graph}
\label{ch:taskA-rozwiazanie}
\label{ch:taskA-ewaluacja}
\label{ch:taskA-benchmark}

Task~A asks whether directed link prediction on a heterogeneous pharmacotherapy graph improves when a learned structural representation is combined with an explicit statistical description of patient-level dependence around each candidate edge.
The structural branch encodes training topology; the statistical branch encodes empirical dependence estimated from training patients; a shared decoder fuses both sources into an edge score.

\section{Problem Formulation}
\label{sec:taskA-problem}

We address directed link prediction among observable nodes of the Synthetic Pharmacotherapy Graph~v3 (Chapter~\ref{ch:dataset}).
For a candidate directed edge $A\to G$, the model produces a score
\begin{equation}
\label{eq:taskA-score}
s_{AG}=f_{\theta}(A,G;G_{\mathrm{train}},X_{\mathrm{train}}),
\end{equation}
where $G_{\mathrm{train}}$ is the message-passing graph available during training and $X_{\mathrm{train}}$ denotes the training-patient observations used to construct node features and pairwise statistical descriptors.
The score is transformed through the sigmoid function,
\begin{equation}
\label{eq:taskA-objective}
\hat{p}_{AG}=\sigma(s_{AG}),
\end{equation}
and the resulting value is used in the binary edge-classification objective.

A positive label indicates that the directed relation is present in the audited graph, whereas a negative label denotes a type-allowed non-edge.
The prediction is based on two complementary sources of evidence:
\begin{itemize}
\item \emph{graph structure}: the typed neighbourhood of $A$ and $G$ represented in $G_{\mathrm{train}}$;
\item \emph{patient-derived statistical evidence}: empirical dependence patterns estimated from training patients for the local motif surrounding the candidate edge.
\end{itemize}
Directionality is determined by the ordered candidate definition and the graph schema, while the pairwise dependence coefficients provide statistical evidence about the variables involved in the candidate relation.

\section{Task-Specific Data Preparation}
\label{sec:taskA-data}

The shared graph and patient scenarios are defined in Chapter~\ref{ch:dataset}.
This section specifies how that resource is turned into prediction instances, a message-passing graph, and node features for Task~A.
The candidate edge inventory is held fixed across observational scenarios, so scenario differences arise from changes in patient-level observability rather than from changes in which edges are scored.

\subsection{Candidate edges and splits}
\label{sec:taskA-candidates}

After restricting the audited graph to observable nodes, the candidate set contains $366$ positive directed edges and $1{,}098$ negatives ($r_{\mathrm{neg}}=3$), for $1{,}464$ candidates in total.
Negatives are drawn from non-edges that satisfy
\begin{equation}
\label{eq:taskA-neg-conditions}
u\neq v,\qquad (u,v)\notin E^{+},\qquad \bigl(\mathrm{type}(u),\mathrm{type}(v)\bigr)\in T_{\mathrm{allowed}},
\end{equation}
rather than from a uniform draw over all ordered node pairs.

Candidates are split independently of the patient split, at edge level, into train/validation/test partitions with proportions $70/15/15$, stratified by the binary edge label.
This yields $256/55/55$ positives and $768/165/165$ negatives.
Patient records are not prediction instances: the training-patient partition is used only to estimate node statistics and pairwise dependence, while labels are attached to candidate edges.

\subsection{Message-passing graph}
\label{sec:taskA-mp-graph}

The HGT encoder operates on a heterogeneous training graph $G_{\mathrm{train}}$ that is smaller than the full audited graph.
Only observable nodes enter the encoder; latent generator nodes are excluded.
Node types and relation names follow the audited v3 ontology (Chapter~\ref{ch:dataset}).

For each positive training relation $u\xrightarrow{r}v$, an auxiliary reverse relation $v\xrightarrow{r^{-1}}u$ is added for message passing.
Reverse relations have distinct identities and enable bidirectional information flow without implying that the underlying mechanistic relation is reversible.

Message passing uses only positive training edges from the candidate set.
Validation positives, test positives, and all negative candidates are excluded from $G_{\mathrm{train}}$.
Held-out candidates are scored by referencing their endpoints in the node embedding table; they are never inserted as edges into $G_{\mathrm{train}}$.

\subsection{Node features}
\label{sec:taskA-node-feats}

Each observable node $v$ receives a three-dimensional empirical feature vector computed from training patients only,
\begin{equation}
\label{eq:taskA-node-x}
x_v=
\bigl[
\mu^{\mathrm{train}}_v,\;
\sigma^{\mathrm{train}}_v,\;
m^{\mathrm{train}}_v
\bigr],
\end{equation}
where $\mu^{\mathrm{train}}_v$ and $\sigma^{\mathrm{train}}_v$ are the train-patient mean and standard deviation of $v$, and $m^{\mathrm{train}}_v$ is the train missingness rate.
Identity embeddings are not used.
Before message passing, each node type is projected independently,
\begin{equation}
\label{eq:taskA-input-proj}
h_v^{(0)}=\mathrm{Dropout}\!\bigl(
\mathrm{ReLU}\bigl(
\mathrm{LayerNorm}(W_{\mathrm{type}(v)}\,x_v)
\bigr)\bigr),
\end{equation}
with $W_{\mathrm{type}(v)}:\mathbb{R}^{3}\to\mathbb{R}^{d}$.
Input statistics are fixed after preprocessing; projection and convolution weights are learned.

\section{Model Architecture}
\label{sec:taskA-arch}

\subsection{Architecture Overview}
\label{sec:taskA-arch-overview}

Figure~\ref{fig:taskA-arch-overview} summarizes the full model.
The structural branch maps $G_{\mathrm{train}}$ through HGT to a graph pair representation $g_{\mathrm{graph}}$.
The statistical branch maps the local $AZ$/$AG$/$ZG$ motif through role-specific encoders to $g_{\mathrm{stat}}$.
Concatenation and a shallow decoder produce the edge logit.

\begin{figure}[htbp]
\centering
\includegraphics[width=\textwidth]{fig_taskA_arch_simplified.png}
\caption{Solution architecture for Task~A.
Patient records yield train-only node features $x_v=[\mu,\sigma,m]$ and pair statistics; candidate edges are split independently at edge level.
The structural branch maps $G_{\mathrm{train}}$ through HGT to $g_{\mathrm{graph}}\in\mathbb{R}^{64}$; the statistical branch encodes roles $AZ$/$AG$/$ZG$ with separate encoders and masks absent roles to $g_{\mathrm{stat}}\in\mathbb{R}^{24}$; fusion yields an edge logit.}
\label{fig:taskA-arch-overview}
\end{figure}

\subsection{Structural Branch}
\label{sec:taskA-hgt}

The graph encoder is a Heterogeneous Graph Transformer~\cite{hu2020hgt} with
\begin{equation}
\label{eq:taskA-hgt-hparams}
d=32,\qquad L=2,\qquad H=4,
\end{equation}
dropout $0.25$, LeakyReLU activations, and a residual connection with LayerNorm after each convolution layer.
These residual links are standard skip connections inside the encoder; they are not studied as an anti-oversmoothing intervention.

After $L$ layers, each observable node has an embedding $z_v\in\mathbb{R}^{32}$.
For candidate $A\to G$, the ordered pair representation is
\begin{equation}
\label{eq:taskA-qag}
q_{AG}=\bigl[z_{A}\parallel z_{G}\parallel z_{A}\odot z_{G}\parallel \lvert z_{A}-z_{G}\rvert\bigr]
\in\mathbb{R}^{128}.
\end{equation}
An MLP graph branch ($128\to 128\to 64$ with GELU and LayerNorm) compresses $q_{AG}$ to
\begin{equation}
g_{\mathrm{graph}}\in\mathbb{R}^{64}.
\end{equation}

\subsection{Statistical Evidence Branch}
\label{sec:taskA-hcr}

The general Hierarchical Correlation Reconstruction (HCR) formulation is stated in Chapter~\ref{ch:related-work}.
Task~A uses an HCR-inspired fixed-width pair vector so that binary, count, continuous and mixed pairs share one downstream interface.
HCR coefficients quantify empirical dependence; they do not identify causal direction.

For an ordered pair $(U,V)$, let $I_{UV}$ be the set of training patients with both variables observed and $n=|I_{UV}|$.
Type-specific orthonormal bases $\{\phi^{(U)}_{j}\}$ and $\{\phi^{(V)}_{k}\}$ yield coefficients
\begin{equation}
\label{eq:taskA-ajk}
a_{jk}
=
\frac{1}{n}
\sum_{i\in I_{UV}}
\phi^{(U)}_{j}(U_i)\,
\phi^{(V)}_{k}(V_i),
\end{equation}
zero-padded to a $4\times 4$ matrix and stored in the first $16$ dimensions of
\begin{equation}
\label{eq:taskA-h40}
h(U,V)\in\mathbb{R}^{40}.
\end{equation}
A binary variable contributes one non-constant contrast; continuous and jittered count variables may contribute up to four low-order basis functions, so the active block ranges from $1\times 1$ (binary--binary) to $4\times 4$.

The remaining dimensions summarize dependence magnitude and complexity, marginal context, joint activity and rarity, estimability, and variable type (Table~\ref{tab:taskA-hcr40}).
All quantities are fitted on training patients only.
An estimability indicator distinguishes unsupported pairs (coefficients zeroed) from genuine near-zero dependence.
Most components are intrinsically bounded or constructed in standardized bases; LayerNorm is applied inside the nonlinear pair encoder.
Complete formula-level definitions of the type-specific bases and all $40$ vector components are provided in Appendix~\ref{app:hcr-details}.

\begin{table}[!ht]
\centering
\footnotesize
\setlength{\tabcolsep}{4pt}
\caption{Grouped composition of the $40$-dimensional HCR-inspired pair representation.}
\label{tab:taskA-hcr40}
\begin{tabular}{@{}lcl@{}}
\toprule
Block & Dim.\ & Purpose \\
\midrule
HCR-inspired coefficients & $16$ & padded dependence-structure matrix $A_{\mathrm{pad}}$ \\
Dependence summaries & $4$ & magnitude and complexity of retained coefficients \\
Marginal context & $8$ & type-dependent descriptors of $U$ and $V$ \\
Joint activity / rarity & $4$ & co-activation enrichment and rarity \\
Support & $2$ & complete-case fraction and estimability flag \\
Variable types & $6$ & one-hot types of $U$ and $V$ \\
\bottomrule
\end{tabular}
\end{table}

For binary--binary pairs, signed Pearson $\phi$ is the natural first-order dependence coefficient and corresponds to the lowest-order HCR contrast; richer blocks retain higher-order structure that a single scalar association cannot express.

\subsection{Fusion and Prediction}
\label{sec:taskA-fusion}

For each candidate $A\to G$, a local context node $Z$ defines three ordered roles $AZ$, $AG$ and $ZG$ (source--context, candidate, context--target).
The same $40$-dimensional construction is applied to each role, yielding a raw motif of $120$ dimensions.
Each role is mapped by an independent MLP encoder
\begin{equation}
\label{eq:taskA-fr}
f_{r}:\mathbb{R}^{40}\to\mathbb{R}^{16}\to\mathbb{R}^{8},
\end{equation}
followed by a hard role gate
\begin{equation}
\label{eq:taskA-role-gate}
e_{r}=m_{r}\,f_{r}(h_{r}),\qquad m_{r}\in\{0,1\},
\end{equation}
so that absent roles (especially missing $Z$) contribute exactly zero.
The statistical branch output is
\begin{equation}
\label{eq:taskA-gstat}
g_{\mathrm{stat}}=[e_{AZ}\parallel e_{AG}\parallel e_{ZG}]\in\mathbb{R}^{24}.
\end{equation}

\begin{figure}[!ht]
\centering
\resizebox{0.62\textwidth}{!}{\begin{tikzpicture}[
  font=\scriptsize,
  >=Stealth,
  every node/.style={align=center},
  pair/.style={inner sep=1pt},
  box/.style={draw, rounded corners=1.5pt, inner sep=2.5pt, minimum height=0.48cm},
  h/.style={box, fill=orange!12, minimum width=1.55cm},
  enc/.style={box, fill=teal!10, minimum width=1.85cm},
  gout/.style={box, fill=green!12, minimum width=3.4cm},
  arr/.style={->, thick, draw=gray!70}
]
% top: ordered pairs
\node[pair] (paz) at (0.0,2.35) {$(A,Z)$};
\node[pair] (pag) at (2.6,2.35) {$(A,G)$};
\node[pair] (pzg) at (5.2,2.35) {$(Z,G)$};

% 40D vectors
\node[h] (haz) at (0.0,1.55) {$h_{AZ}$ 40D};
\node[h] (hag) at (2.6,1.55) {$h_{AG}$ 40D};
\node[h] (hzg) at (5.2,1.55) {$h_{ZG}$ 40D};

% encoder + mask
\node[enc] (eaz) at (0.0,0.75) {encoder + mask};
\node[enc] (eag) at (2.6,0.75) {encoder + mask};
\node[enc] (ezg) at (5.2,0.75) {encoder + mask};

% concat + g_stat
\node[gout] (cat) at (2.6,-0.05) {$[e_{AZ}\,\Vert\,e_{AG}\,\Vert\,e_{ZG}]$};
\node[gout] (gst) at (2.6,-0.85) {$g_{\mathrm{stat}}\in\mathbb{R}^{24}$};

\draw[arr] (paz) -- (haz);
\draw[arr] (pag) -- (hag);
\draw[arr] (pzg) -- (hzg);
\draw[arr] (haz) -- (eaz);
\draw[arr] (hag) -- (eag);
\draw[arr] (hzg) -- (ezg);
\draw[arr] (eaz.south) |- (cat.west);
\draw[arr] (eag) -- (cat);
\draw[arr] (ezg.south) |- (cat.east);
\draw[arr] (cat) -- (gst);
\end{tikzpicture}
}
\caption{Role-aware motif for $A\to G$: three ordered pair roles, each with a separate encoder and role mask, concatenated to $g_{\mathrm{stat}}\in\mathbb{R}^{24}$.}
\label{fig:taskA-motif}
\end{figure}

Fusion concatenates the branches,
\begin{equation}
g_{\mathrm{fusion}}=\bigl[g_{\mathrm{graph}}\parallel g_{\mathrm{stat}}\bigr]\in\mathbb{R}^{88},
\end{equation}
and maps $88\to 64\to 1$ (GELU, LayerNorm, dropout) to a single edge logit.

\section{Training and Evaluation}
\label{sec:taskA-train-eval}

\subsection{Training Objective}
\label{sec:taskA-loss}

For each candidate $i$ with logit $s_i$ and label $y_i\in\{0,1\}$,
\begin{equation}
\label{eq:taskA-sigmoid}
\hat{p}_{i}=\sigma(s_{i}),
\end{equation}
and parameters are optimized by weighted binary cross-entropy,
\begin{equation}
\label{eq:taskA-bce}
\mathcal{L}
=
-\frac{1}{N}
\sum_{i=1}^{N}
\bigl[
w_{+}\, y_{i}\log\hat{p}_{i}
+(1-y_{i})\log(1-\hat{p}_{i})
\bigr],
\end{equation}
with $w_{+}=N_{\mathrm{neg}}/N_{\mathrm{pos}}\approx 3$ under the $3{:}1$ negative-sampling ratio.
Weighting balances class contributions to the loss; it does not change the candidate inventory.

\subsection{Evaluation Metrics}
\label{sec:taskA-ap}

Model selection uses validation Average Precision (AP)~\cite{sklearn,davis2006},
\begin{equation}
\label{eq:taskA-ap}
\mathrm{AP}
=
\sum_{k}
\bigl(R_{k}-R_{k-1}\bigr)\,P_{k},
\end{equation}
where $P_k$ and $R_k$ are precision and recall at successive thresholds on the ranked score list ($R_0=0$).
This is the Average Precision summary of the ranked prediction list used throughout the thesis.
AP is appropriate for imbalanced link prediction because it emphasizes precision as recall increases.
The sealed test set is never used to choose architecture, representation or hyperparameters.
When discrete metrics such as $F_1$ are reported, the classification threshold is selected on validation only.

\subsection{Model and Training Configuration}
\label{sec:taskA-train-config}

Table~\ref{tab:taskA-train} lists the model and training configuration of the full model.
Six observational scenarios are evaluated; the official macro averages five scenarios and excludes the hidden-confounder scenario, which is reported for completeness only.

\begin{table}[!ht]
\centering
\footnotesize
\setlength{\tabcolsep}{4pt}
\caption{Task~A: model and training configuration.}
\label{tab:taskA-train}
\begin{tabular}{ll}
\toprule
Setting & Value \\
\midrule
Optimizer & Adam~\cite{kingma2015adam} \\
Learning rate & $10^{-3}$ \\
Weight decay & $5\times 10^{-4}$ \\
Max epochs & $200$ \\
Early stopping & patience $40$ on validation AP \\
Gradient clip & $1.0$ (global norm) \\
Normalization & LayerNorm in HGT and decoder blocks~\cite{ba2016layernorm} \\
HGT & $d=32$, $L=2$, $H=4$, dropout $0.25$~\cite{hu2020hgt} \\
Statistical branch & HCR-inspired $40$D pair vectors; MLP role encoders \\
Fusion & $64+24=88$ dimensional late fusion \\
Training seeds & five fixed seeds \\
Selection metric & validation AP \\
Official macro & mean over five scenarios (excludes hidden confounder) \\
\bottomrule
\end{tabular}
\end{table}

\section{Experimental Results}
\label{sec:taskA-results}

Results address three questions: which graph backbone to use, how much explicit statistical context contributes, and how the full model performs across observational scenarios.
Architecture and hyperparameters were selected on validation AP before sealed-test reporting.

\subsection{Architecture Comparison}
\label{sec:taskA-stage-a}

With the statistical branch disabled and a common decoder, HGT~\cite{hu2020hgt}, R-GCN~\cite{schlichtkrull2018rgcn}, GraphSAGE~\cite{hamilton2017graphsage} and GATv2~\cite{brody2022gatv2} were compared on the clean scenario.
Mean validation AP was approximately $0.729$ (HGT), $0.697$ (R-GCN), $0.671$ (GraphSAGE) and $0.650$ (GATv2).
HGT was retained as the structural encoder for all subsequent experiments.

\subsection{Contribution of Statistical Evidence}
\label{sec:taskA-stage-c}

With HGT fixed, explicit pairwise statistics were enriched progressively on the clean scenario (three seeds).
Validation AP rose from approximately $0.719$ with no pairwise statistics, to $0.840$ with signed Pearson $\phi$, $0.853$ with normalized mutual information, $0.913$ with an enriched multi-component representation, and $0.917$ with the full $40$-dimensional HCR-inspired vector.
The graph-only value $\approx 0.719$ in this ablation is not identical to the $\approx 0.729$ reported in the backbone comparison above: the former is the HGT-only arm of the statistical-enrichment ladder under a shared decoder and seed set, whereas the latter compares alternative graph encoders under a matched backbone-selection protocol.
Across the five official scenarios the full representation reached a macro validation AP of approximately $0.915$.
Removing the four derived dependence-summary dimensions reduced validation AP by about $0.008$ on average; the full vector was retained.

\subsection{Final Evaluation}
\label{sec:taskA-final}

Figure~\ref{fig:taskA-final-metrics} reports the full model under the official multi-scenario protocol.
Official macro test AP is $0.894$ and test AUROC is $0.967$.
Performance remains stable under selection bias, noisy documentation and related observational artefacts; the largest train--test AP gap appears in the multi-hospital scenario.

\begin{figure}[htbp]
\centering
\includegraphics[width=\textwidth]{fig_final_metrics_per_scenario.pdf}
\caption{Task~A: final performance across observational scenarios.
Panel~(a): AP on the training partition (at the validation-selected epoch), on validation, and on the sealed test set.
Panel~(b): AUROC on validation and sealed test.
Bars are means over five training seeds; whiskers are one standard deviation.
The hatched scenario is excluded from official macro averages.
Official macro test AP is $0.894$ and test AUROC is $0.967$ (five-scenario mean; hidden confounder excluded).}
\label{fig:taskA-final-metrics}
\end{figure}
